(a) A line cuts out a divisor linearly equivalent to $3P_0$, so collinearity, counted with intersection multiplicities, implies $P+Q+R=0$. Conversely, let the line through $P,Q$ meet $X$ a third time at $R'$; use the tangent if $P=Q$. Then $P+Q+R'=0$, whence $R'=R$.

(b) The tangent at $P$ cuts out $2P+R$, where $2P+R=0$. It passes through $P_0$ exactly when $2P=0$. Thus for $P\neq P_0$ this is equivalent to order two; $P_0$ itself has order one and also satisfies the tangent condition.

(c) The tangent has intersection multiplicity three at $P$ exactly when its third intersection is $P$. By (a), this is equivalent to $3P=0$. Thus the inflection points other than $P_0$ have order three, while $P_0$ is also an inflection point.

(d) The chord and tangent constructions use rational coefficients for rational points, and the remaining intersection is rational. Also $-(x,y)=(x,-y)$ and $P_0\in X(\mathbb Q)$, so $X(\mathbb Q)$ is a subgroup.

If $(x,y)\in X(\mathbb Q)$ and $y\neq0$, then the positive numbers $|(x^2-1)/y|$, $|2x/y|$, and $|(x^2+1)/y|$ are the sides of a rational right triangle of area one. We show that no such triangle exists. Clearing denominators and choosing a counterexample with least integral hypotenuse gives a primitive integer right triangle with square area. Its sides are $2mn,m^2-n^2,m^2+n^2$, with $m>n$ coprime and of opposite parity. The four pairwise coprime integers $m,n,m+n,m-n$ have square product, so each is a square. Write $m+n=a^2$ and $m-n=b^2$. Then $a,b$ are odd, and
\[\left(\frac{a+b}{2}\right)^2+\left(\frac{a-b}{2}\right)^2=m.\]
This is another integral right triangle, with hypotenuse $\sqrt m$ and area $n/4$. Since $a^2-b^2=2n$ is divisible by $8$ and $n$ is a square, $n/4$ is a square. Its hypotenuse is smaller, a contradiction.

Consequently $y=0$, so the rational points are exactly $P_0,(0,0),(1,0),(-1,0)$. Each of the last three has order two, and hence $X(\mathbb Q)\cong(\mathbb Z/2\mathbb Z)^2$.
